\begin{frame}{Solución angular}
	Trabajando sobre la parte angular de las ecuaciones \ref{eq:separacion de variables RY}
	\[
		\frac{1}{Y}\left[ \frac{1}{\sin{\theta}} \frac{\partial}{\partial\theta} \left(\sin{\theta}\frac{\partial Y}{\partial\theta}\right)+ \frac{1}{r^2\sin^2\theta} \frac{\partial^2Y} {\partial\phi^2}\right] =-\lambda
	\]
	Plantearemos una nueva separación de variables
	\begin{equation}
		Y(\theta,\phi) = \Theta(\theta) \Phi(\phi)
		\label{eq:propuesta separacion de variables angulares}
	\end{equation}
\end{frame}

\begin{frame}{Solución angular}
	Por lo que nos queda:
	\begin{equation}
		\begin{cases}
			\begin{aligned}
			   \frac{\sin{\theta}}{\Theta}\frac{d}{d\theta} \left( \sin{\theta} \frac{d\Theta}{d\theta}\right)+\lambda \sin^2{\theta} &= m^2 \\
				\frac{1}{\Phi}\frac{d^2\Phi}{d\phi^2} &=-m^2
			\end{aligned}
		\end{cases}
    \label{eq:separacion de variables angular}
	\end{equation}
\end{frame}
